
\section{Testable questions on the supersymmetry of the schemes}
\label{sec:questions}

In this formulation the supersymmetry of a dimensional scheme at a given loop order is the statement $\mathrm{Ob}=0$ in the components along $\xi$, $\eta$ and their arity-two closure, without finite counterterms; by Theorem \ref{thm:qap} it is a property of the insertions of the evanescent breaking $\widehat\Delta$ alone [Bar98, §2.8]. High-loop computations of $\mathcal N=4$ data in dimensional schemes [FSSZ08], [Vel08], [KLOV04] assume this property at their loop order; no Slavnov--Taylor or supersymmetry Ward identity for composite operators has been evaluated in DRED, superspace DR or BMHV at any loop order [LIT, §2.3]. The questions below are those of [LIT, §5], numbered Q1--Q9, in the notation of this note; what the computation must reproduce is Definition \ref{def:targets}.

\begin{proposition}[No evanescent couplings in Scheme D]
\label{prop:noevcoup}
The rotations of the internal space $E=Q(4{-}D)S\oplus W^\perp_{\mathbb R}$, acting on $A|_E$ as a vector and on $\lambda$ through $\rho\circ q$, are an exact symmetry of $S_D$, $\Psi$ and the non-minimal sector. Hence every counterterm of Scheme D is $\mathfrak{so}(E)$-covariant: the $4-D$ $\epsilon$-scalars and the six scalars renormalise with one wave-function constant, one quartic and one Yukawa structure, all those of the reduction of ten-dimensional $F_A^2$ and $(\lambda,{\not d}_A\lambda)$; the separate evanescent couplings of dimensional reduction for non-supersymmetric gauge theories have no counterpart.
\end{proposition}

\begin{proof}
$S_D$ is the reduction of the ten-dimensional action to $QDS$; rotations of the reduced directions commute with the derivatives along $QDS$ and preserve the Hodge pairing, the Clifford relation and $(\rho(v)\zeta_1,\zeta_2)=(\zeta_1,\rho(v)\zeta_2)$; $\Psi$ contains $A|_{QDS}$ only. The counterterms are local functionals built from the vertices and propagators, all $\mathfrak{so}(E)$-covariant, so the pole parts are $\mathfrak{so}(E)$-invariant; the invariant monomials of the relevant weight are those listed.
\end{proof}

\begin{problem}[Q1: composite operators, Scheme D, one loop]
\label{q:Q1}
For a gauge-invariant $O$ let $A_O$ be the one-loop part of $L(O)$ along $\xi$ in Scheme D with $\psi=0$ and $F_1=0$ along $\xi$ (Theorem \ref{thm:master}: the Fierz vertex $\mathfrak F(\xi)$ on the cycles through two letters, the counterterm commutator $\mathcal C_\xi$, the operator breaking). Prerequisite: $\mathrm{Ob}^{(1)}(\xi)=0$ and $\mathrm{Ob}^{(2)}(\xi,\xi')=0$ at the level of the action (gate G4c; the ten-dimensional form of the counting rules [S05, (4.10)] decides it). Decide:
\begin{itemize}
\item[(a)] $A_O=0$ on the stress-tensor multiplet and the other half-BPS multiplets, in particular on the $\mathbf{20}'$ bilinears. By (I5) $A_O$ is a coboundary $[\rho_0,e]$ there in every scheme; the question is whether Scheme D needs $e\neq0$, i.e. a finite operator counterterm.
\item[(b)] On the Konishi multiplet, $A_O=[\rho_0,e]+$(the lifting): by (I4) and Remark \ref{rem:twistDR} the lifting $\kappa$ is non-zero and obeys $\kappa(\eta)\beta(\xi)c_0=-\gamma_K(\xi,\eta)c_0$; it is not removable by finite counterterms. The question is the remainder: is it a coboundary, and with which $e$?
\item[(c)] $e=0$ (no finite restoring operator counterterm, as for elementary fields at one loop in SQED and SQCD [HKS99], [HS01]) or $e\neq0$. Letters $\bar\lambda\,\rho\cdots\rho\,\lambda$ lengthen the open $\rho$-strings through the operator and are the candidates for $e\neq0$; $e\neq0$ shifts the finite part of $Z_O$ and two-loop mixing coefficients such as the Konishi-anomaly coefficient [BKRS01].
\end{itemize}
Loop order one; Route A with Route B as check; gates G4c, G14.
\end{problem}

\begin{problem}[Q2: Scheme B, one loop]
\label{q:Q2}
\begin{itemize}
\item[(a)] $\mathrm{Ob}\neq0$ already on elementary functions: compute the $\mathcal N=4$ analogues of the one-loop SQED and SQCD restoring counterterms [HKS99], [HS01] (gauge--Yukawa and gauge--quartic relations) as the components of $F_1^{\rm B}$ (Theorem \ref{thm:restore}, gate G4d).
\item[(b)] The finite operator counterterms $e^{\rm B}$ for the $\mathbf{20}'$ and Konishi multiplets; before the $R$-restoration the $SO(4)_R$-singlet component of the $\mathbf{20}'$ and the Konishi bilinear share all exact quantum numbers and mix.
\item[(c)] Two loops: with $F_1^{\rm B}$ and $e^{\rm B}$ the two-loop dilatation operator equals that of Scheme D up to a similarity transformation (Theorem \ref{thm:D1}, [DG91], [HN95]); without them the $\mathbf{20}'$ acquires $\gamma^{(2)}\neq0$.
\item[(d)] Is $F_1^{\rm B}$, up to $s$-exact terms, the finite counterterm produced by the missing $4-D$ $\epsilon$-scalars, $\mathrm{LIM}$ of $(4-D)\times$ the residues of the $\epsilon$-scalar cycles of Scheme D? This identifies the one-loop conversion between the schemes as a canonical transformation (Proposition \ref{prop:lawsD} (f)).
\item[(e)] With background $SU(4)_R$ gauge fields (position-dependent $R$-ghosts) the only non-exact class must be the 't Hooft anomaly, proportional to $(N^2-1)d^{abc}$; with constant ghosts it does not occur.
\end{itemize}
Loop order one for (a), (b), (d), (e); two for (c).
\end{problem}

\begin{problem}[Q3: superconformal closure on operators, one loop]
\label{q:Q3}
Relation (R11): on the $\mathfrak{su}(2|3)$ sector, $c_9$ and $c_{10}$ are the coefficients of the length-changing parts of the classical supercharges [Bei03b, (3.32)], read from the tree-level BV data of Proposition \ref{prop:QS} and Definition \ref{def:equiv}, $c_1$ is the normalisation of the one-loop dilatation operator from Theorem \ref{thm:D1}, and the order-$g^2$ closure of $\{Q,S\}$ requires $c_9c_{10}=\frac12c_1$ [Bei03b, (3.36)], the $\mathfrak{su}(2|3)$ form of $\kappa(\eta)\beta(\xi)c_0=-\gamma(\xi,\eta)c_0$ [FR, Thm 4.9, §7]. Decide whether it holds in Scheme D with $F=0$, and in Scheme B only with $F_1^{\rm B}$, $e^{\rm B}$. This turns the identification of the deformed $\mathfrak{su}(2|3)$ representation with renormalised operators, assumed in [Bei03b], into a computation. Loop order one; gate G13.
\end{problem}

\begin{problem}[Q4: conformal breaking at $D\neq4$, one loop]
\label{q:Q4}
$\mathcal N=4$ in $D$ dimensions has no interacting fixed point; conformal symmetry is broken by $(4-D)\int\sigma_v\mathscr L_0$ (Theorem \ref{thm:breakD} (ii), [DHKS07, (25), (27)]).
\begin{itemize}
\item[(a)] Normalisation: Theorem \ref{thm:D1}, $D_1=-2\,\mathrm{LIM}\,\mathcal R$, against the Konishi and $\mathbf{20}'$ values (gates G5, G6).
\item[(b)] $K_1^0$ of Theorem \ref{thm:K1} is the one-loop special conformal anomaly; at two loops the order-$\hbar^2$ part of the $(K,P)$-component of $L\circ L=0$, $[K^{(0)},D^{(2)}]=[D^{(1)},K^{(1)}]$, fixes the non-conformal part of the two-loop dilatation operator [BKM05].
\item[(c)] In Scheme B, does $K_1$ acquire the extra term $[\rho_0(K),e^{\rm B}]$ from the supersymmetry-restoring operator counterterm, i.e. do conformal covariance and supersymmetry restoration require the same finite counterterm?
\end{itemize}
Loop order one for (a), (c); two for the constraint in (b).
\end{problem}

\begin{problem}[Q5: the four-scalar identity, two loops]
\label{q:Q5}
The identity tying the quartic scalar vertex to $g^2$ and to the Yukawa vertex is established in dimensional reduction at one loop [S05, (4.8)] and, in a restricted limit, at three loops [LIT, §2.2]. In Scheme D: is the two-loop Fierz-insertion part of $\mathrm{Ob}^{(2)}$ on the four-scalar function zero? Proposition \ref{prop:noevcoup} fixes the pole structure ($\epsilon$-scalars and scalars renormalise together); the question is the finite part. Loop order two.
\end{problem}

\begin{problem}[Q6: onset of the breaking of dimensional reduction, four and five loops]
\label{q:Q6}
Determine the ten-dimensional counting rule: the minimal length of a closed $\rho$-string in $Q4S\oplus W^\perp$ for which an insertion of $\mathfrak F$ survives in the formal algebra (the four-dimensional rule gives the five-index structure with trace $48D(D-1)(D-2)(D-3)(D-4)$ and four loops for the pure-SYM propagator [LIT, §2.4]). Test the claims that the insertion vanishes on the $\mathcal N=4$ Yukawa vertex through four loops, with the first candidates at five loops (six flavour matrices in one trace), and settle the conflicting estimates of five versus eight or ten loops [LIT, §2.4]. Include the chirality-odd traces: by Remark \ref{rem:g11} they enter no pole at one loop; find the first order at which a closed fermion loop with at least ten $\rho$'s reaches a pole. Loop orders four and five, by counting and selected graphs.
\end{problem}

\begin{problem}[Q7: local coupling, one loop]
\label{q:Q7}
Promote the scale of the pairing of $\mathfrak g$ to a background function. The breakings $\widehat\Delta^{(1)}(\Delta)$, $\widehat\Delta^{(1)}(K_{v_1})$ are then the variations along $\sigma_v$-weighted coupling directions. Compute the one-loop local-coupling class of the constant-ghost Slavnov--Taylor identity in Scheme D (in the $\mathcal N=1$ language it is $r_\eta=C(G)/8\pi^2$ for $\mathcal N=1$ SYM [Kra02] and absent for $\mathcal N=2$ [LIT, §3]); its vanishing for $\mathcal N=4$ would make the Lagrangian-insertion formula $\partial_{g^2}\langle\cdots\rangle=\int\langle\mathscr L\cdots\rangle$ exact up to contact terms. Loop order one.
\end{problem}

\begin{problem}[Q8: deformed theories, one loop]
\label{q:Q8}
For $\beta$- and $\gamma_i$-deformed $\mathcal N=4$ the local cochains must include double-trace couplings; their one-loop beta functions are the obstruction to the conformal components of the master equation; $U(N)$ and $SU(N)$ differ by the abelian factor. The construction of Sections \ref{sec:schemes}--\ref{sec:oneloop} applies with the deformed vertices; not implemented in the engines of Section \ref{sec:engines}. Loop order one.
\end{problem}

\begin{problem}[Q9: $\gamma_5$ in the Standard Model]
\label{q:Q9}
The BV formulation of Scheme B with sterile partners (Proposition \ref{prop:nogo}, the genuinely chiral case) is the setting of [BIMS20], [BIKMS21]. Statement to test: with symmetry-restoring counterterms the $\gamma_5$-odd four-loop coefficient of the strong beta function equals the reading-point value fixed by Weyl consistency conditions [LIT, §4]; smallest test, a chiral gauge--Yukawa model at the order where two $\gamma_5$-odd traces first meet in a divergent graph. Outside the present engines; loop order four in the Standard Model.
\end{problem}

\begin{problem}[Classes across the tracks]
\label{q:classes}
By Theorem \ref{thm:inv} the data (I1)--(I6) agree between Scheme D, Scheme B and the heat-kernel track once the latter is built on a valid gauge fixing. A difference would show that (Hs) fails, i.e. an additional closed direction of the one-loop quantisation. Gate G12.
\end{problem}

\begin{remark}[What is settled where]
Q1, Q2 (a), (b), (d), Q3, Q4 (a), (c), Q7 and the class comparison are one-loop questions settled by Route A with Route B as the independent check; Q2 (c), Q4 (b), Q5 require the two-loop extension of Theorem \ref{thm:master} (three letters at most, [HK, DESIGN\_STEP2, R8]) and of Lemma \ref{lem:bonneau} [BIKMS21]; Q6 is a counting statement followed by selected graphs; Q8, Q9 apply the framework outside the undeformed $\mathcal N=4$ theory.
\end{remark}
